\documentclass[12pt]{article}
\usepackage{amsmath, amssymb, geometry, graphicx, tikz}
\usepackage{actuarialangle}
\usetikzlibrary{patterns}
\usetikzlibrary{calc}
\geometry{margin=1in}
\setlength{\parskip}{1em}
\setlength{\parindent}{0pt}

\begin{document}

\begin{center}
    \Large \textbf{Time value of money formulas}
\end{center}

(1.) Amount of Interest from $t_1$ to $t_2$ if \$$K$ is invested: $A_K(t_2)-A_K(t_1)$

(1.) Effective interest rate from $t_1$ to $t_2$: $i_{[t_1,t_2]} = \dfrac{a(t_2)-a(t_1)}{a(t_1)}$

(2.) Effective discount rate from $t_1$ to $t_2$:  $d_{[t_1,t_2]} = \dfrac{a(t_2)-a(t_1)}{a(t_2)}$

(3.) Simple Interest accumulation function: $a(t) = 1+rt$ (where $r = i_{[0,1]}$)

(4.) Simple Discount accumulation function: $a(t) = \frac{1}{1-rt}$ (where $r = d_{[0,1]}$)

(4.) Compound Interest accumulation function: $a(t) = (1+i)^t$ (where $i = i_{[t,t+1]}$ for any $t$)

(5.) Discount function: $v(t)=\dfrac{1}{a(t)}$ and for compound interest $v = v(1)= \dfrac{1}{1+i}$

(6.) $1+i = \left(1+\frac{i^{(m)}}{m}\right)^m $ where $i^{(m)}$ is the nominal interest rate compounded $m$ times per year

(7.) $1-d = \left(1-\frac{d^{(m)}}{m}\right)^m $ where $d^{(m)}$ is the nominal discount rate compounded $m$ times per year

(8.) $d = \dfrac{i}{1+i}$ and $i = \dfrac{d}{1-d} = \dfrac{d}{v}$

(9.) Constant force of interest or nominal rate convertible continuously: 
\vspace{-.5cm}
    
    \[\delta = \ln(1+i) \iff e^{\delta}=1+i\]


(10.) Force of interest at time $t$: $\delta_t = \dfrac{a'(t)}{a(t)} = \dfrac{d}{dt}\left(\ln(a(t)\right) \iff a(t) = e^{\displaystyle \int_0^t \delta_r \, dr}$




\end{document}
